\documentclass{leveldoc}
\DocUses{foundations,high_voltage_dome}
\DocFingerprint{aae338f3ab90f785}

\begin{document}

\section{What the level asks}\label{sec:asks}

\leveldescription

One positive ion is launched along $+x$ from the left edge. A metal sphere sits
above the line of flight and is isolated, with net charge zero. A strong positive
charge already stands above the sphere, so the metal polarises: the near side
takes the opposite sign and the far side the same sign, while the net charge stays
zero. The player may add up to three charges, of either sign, from the listed
magnitudes, inside the allowed region. There are no magnets.

The ion must enter the detector. There is no acceptance window, no kinetic window
and no gate. A flight ends when the ion touches a charge or the sphere, leaves the
world, or enters the detector (\xsecref{foundations}{sec:board}). The sphere's
obstacle is the metal plus a contact shell of thickness $\val{contact}$.

The reference adds one negative charge just under the line of flight. Without it
the ion is driven down and leaves through the bottom of the world. With it the ion
dips, then climbs into the detector through the detector's left edge. This
document uses the dome's image for one sphere and adds the centre charge that
keeps the sphere neutral, then integrates the flight by the classical
Runge--Kutta method.

\section{Setup}\label{sec:setup}

\setuptable

The speed of light is $c = \val{cLight}$. The ion has charge $q = \val{qPart}$, mass
$m = \val{mPart}$ and radius $0$, and it is launched with kinetic energy
$T_0 = \val{T0}$. A fixed charge has radius $R = \val{Rcharge}$. The sphere has
centre $(\val{sphereX}, \val{sphereY})$, radius $a = \val{sphereA}$ and net charge
$Q = \val{sphereQ}$. The level's charge is $Q_\mathrm{u} = \val{QUp}$ at
$(\val{xUp}, \val{yUp})$. The reference charge is $Q_\mathrm{p} = \val{QPlay}$ at
$(\val{xPlay}, \val{yPlay})$.

The launch is at $(\val{xLaunch}, \val{yLaunch})$, straight along $+x$. Then
$\gamma_0 = \val{gamma0}$, $\abs{\vect{p}_0} = \val{p0}$ and
$v_0 = \val{v0} = \val{vOverC}\,c$ (\xeqref{foundations}{eq:launch}). The detector
is $\val{xDetLo} \le x \le \val{xDetHi}$, $\val{yDetLo} \le y \le \val{yDetHi}$.
The world is $-1 \le x \le \val{xWorld}$, $-1 \le y \le \val{yWorld}$.

\begin{figure}[ht]
\centering
\begin{tikzpicture}
\begin{axis}[axis equal image, width=0.92\linewidth,
  xmin=-1.5, xmax=32.5, ymin=-2, ymax=22.5,
  xlabel={$x$}, ylabel={$y$}, tick label style={font=\small}]
  \draw[black!40, dashed] (-1,-1) rectangle (\rawval{xWorld},\rawval{yWorld});
  \draw[black!25] (0,0) rectangle (30,20);
  \fill[green!50!black, fill opacity=0.18]
    (\rawval{xDetLo},\rawval{yDetLo}) rectangle (\rawval{xDetHi},\rawval{yDetHi});
  \draw[green!40!black]
    (\rawval{xDetLo},\rawval{yDetLo}) rectangle (\rawval{xDetHi},\rawval{yDetHi});
  \fill[blue!30, fill opacity=0.35] (\rawval{sphereX},\rawval{sphereY}) circle[radius=\rawval{sphereA}];
  \draw[blue!50!black] (\rawval{sphereX},\rawval{sphereY}) circle[radius=\rawval{sphereA}];
  \fill[red!70!black] (\rawval{xUp},\rawval{yUp}) circle[radius=\rawval{Rcharge}];
  \fill[purple!70!black] (\rawval{xPlay},\rawval{yPlay}) circle[radius=\rawval{Rcharge}];
  \addplot[orange!80!black] table {pathbare.dat};
  \addplot[very thick, blue!60!black] table {pathref.dat};
  \draw[-{Stealth[length=5pt]}, thick] (\rawval{xLaunch},\rawval{yLaunch}) -- (\rawval{arrowX},\rawval{arrowY});
  \fill (\rawval{xLaunch},\rawval{yLaunch}) circle[radius=0.08];
\end{axis}
\end{tikzpicture}
\caption{The neutral sphere and the two flights. The circle is the sphere, the red
disc is the level's charge and the violet disc is the reference charge; the green
box is the detector. Blue is the reference flight: it dips below the launch line,
then climbs and meets the left edge of the detector. Orange is the same launch
with the reference charge removed: it leaves through the bottom edge of the world.
Both paths stay outside the sphere.}
\label{fig:board}
\end{figure}

\section{Principles and model}\label{sec:model}

\begin{principles}
  \item The field of a point charge and the force on a non-polarizable particle,
    \xeqref{foundations}{eq:coulomb} and \xeqref{foundations}{eq:sphereforce}.
  \item The relativistic equation of motion \xeqref{foundations}{eq:eom}, the launch
    \xeqref{foundations}{eq:launch}, and energy in a static field
    \xeqref{foundations}{eq:energy}. The motion stays in the plane
    (\xsecref{foundations}{sec:plane}).
  \item A metal sphere is an equipotential. Its relaxation is instantaneous on the
    flight's time scale (\physref{2.6}). One sphere is the image of the dome,
    \xeqref{high_voltage_dome}{eq:kelvin}, plus the centre charge of
    \S\ref{sec:neutral}.
  \item The board's rules (\xsecref{foundations}{sec:board}): obstacles first, then
    the world, then the detector. The classical Runge--Kutta method and the cubic
    Hermite location of an event (\xsecref{foundations}{sec:numerics}).
  \item The radiated power \xeqref{foundations}{eq:larmor}, summed by Simpson's rule
    \xeqref{foundations}{eq:simpson}. It is not put back into the force.
\end{principles}

Outside the sphere the electric field is the Coulomb field of the real charges, of
their images and of one charge $Q_c$ at the centre, plus the image of the ion and
the centre charge that follows the ion. There is no magnetic field. The state
$(\vect{x}, \vect{p})$ obeys \xeqref{foundations}{eq:eom} with that $\vect{E}$,
from the launch until the first ending event.

\begin{assumption}{Electrostatic images}\label{as:static}
The metal answers at once. That is exact at infinite $c$, and for finite $c$ it
drops the retardation across the board. The launch speed is $\val{vOverC}\,c$.
The game uses the same static images (\physref{2.6}), so this is the model being
compared, not a correction on top of it.
\end{assumption}

\begin{assumption}{No radiation reaction}\label{as:radiation}
The ion radiates \xeqref{foundations}{eq:larmor}, and that energy is not taken
from the flight. The larger fraction of the launch energy on the two flights here
is $\val{radFracMax}$.
\end{assumption}

\begin{assumption}{The ion's own image}\label{as:self}
The ion induces the pair in \S\ref{sec:self}. Along these two flights the largest
sampled $\abs{q\vect{E}_\mathrm{self}}$ is $\val{selfForce}$, a fraction
$\val{selfBound}$ of $T_0$ per cell. The samples are the step ends and nine
Hermite stations on each step, not a continuous maximum. The force is kept in
the integration. The smallest clearance of the surface is $\val{clearMin}$,
against a shell of thickness $\val{contact}$ and a sample span $\val{sampleArc}$.
The closest centre distance on the two flights is $\val{rNear}$.
\end{assumption}

\begin{assumption}{The fixed step}\label{as:step}
The reported step is $h = \val{H}$. Halving it estimates the largest error of an
end state as $\val{errMax}$ (\xsecref{foundations}{sec:numerics}). Ten times that
estimate lies above $10^{-9}$, so the tolerance on a state is that product,
$\val{tolFlight}$.
\end{assumption}

\section{Derivation}\label{sec:derivation}

\subsection{A neutral isolated sphere}\label{sec:neutral}

The image of one charge in a grounded sphere is
\xeqref{high_voltage_dome}{eq:kelvin}. With one sphere there is nowhere further
to image that charge, so the construction is exact (\physref{2.6}). The images of
the fixed charges sum to a net induced charge $g$. Each image lies inside the
metal, so $g$ is part of the sphere's charge.

The dome holds its potential with a centre charge (\xeqref{high_voltage_dome}{eq:centre}):
a charge at the centre adds its own value to the net charge and a constant to the
potential, and it does not disturb the grounded cancellation. An isolated sphere
has no supply. Its net charge is the prescribed $Q$, so the centre charge and the
potential it produces are
\begin{equation}\label{eq:neutral}
  Q_c = Q - g, \qquad \alpha = Q_c / a.
\end{equation}
Here $Q = \val{sphereQ}$ and $g = \val{inducedG}$, so $Q_c = \val{centreQ}$ and
$\alpha = \val{alpha}$. The images and this centre charge sum to $\val{netQ}$, a
fraction $\val{netRel}$ of $\abs{Q_c}$: roundoff of the cancellation, not a leftover
charge.

The level's charge lies at distance $\val{sourceRUp}$ from the centre, so its image
is $\val{imageQUp}$ at $(\val{imageXUp}, \val{imageYUp})$. The reference charge
lies at distance $\val{sourceRPlay}$, so its image is $\val{imageQPlay}$ at
$(\val{imageXPlay}, \val{imageYPlay})$. On sixteen points of the surface the
potential of this fixed set differs from $\alpha$ by at most $\val{surfRes}$, a
fraction $\val{surfRel}$ of $\abs{\alpha}$.

\subsection{The ion's own image}\label{sec:self}

The ion at distance $r$ from the centre induces the same grounded image. On the
dome the supply absorbs the change of net charge, so no second centre charge
follows the ion. Here the net charge must stay $Q$, and the image charge is
$-qa/r$, so a centre charge $+qa/r$ follows the ion. The two sum to zero. At
launch they are $\val{selfQ}$ and $\val{selfCentre}$, and their sum is
$\val{ownNet}$.

The potential of the image at the ion is \xeqref{high_voltage_dome}{eq:self}. The
following centre charge sits at distance $r$, so its potential at the ion is
$qa/r^2$. The self potential is the sum
\begin{equation}\label{eq:selfphi}
  \phi_\mathrm{self}
  = -\frac{qa}{r^2 - a^2} + \frac{qa}{r^2}.
\end{equation}
Both terms are linear in $q$. The conserved energy is the dome's
\xeqref{high_voltage_dome}{eq:image-energy}, with this $\phi_\mathrm{self}$: the
external potential of the fixed charges, their images and $Q_c$, plus half the
self potential. At launch, $q\phi_\mathrm{ext} = \val{qPhi0}$ and the self part
is $\val{halfSelf0}$.

The force on the ion is $q\vect{E}$ of every charge except the ion, including both
pieces of \eqref{eq:selfphi}. The radial formula that counts only the grounded
image does not apply: the following centre charge pushes back.

\subsection{The flight}\label{sec:flight}

There is no closed orbit. The fixed set is five charges, and the ion's pair moves.
The state is advanced by \xeqref{foundations}{eq:rk4} with the fixed step
$h = \val{H}$. After each step the cubic Hermite interpolant of the two ends is
sampled at sixty-four fractions of the step. The first sample that has touched an
obstacle, left the world or entered the detector is refined by bisection, eighty
halvings (\xsecref{foundations}{sec:numerics}). Obstacles are tested before the
world, and the world before the detector. A face coordinate is assigned when the
residual is below $10^{-6}$.

The error of the fine step is $\abs{y_h - y_{2h}}/15$. The same rule estimates
the Simpson error of the radiated energy. Both flights give the same outcome at
$h$ and at $2h$.

\section{The reference solution by hand}\label{sec:hand}

\begin{procedure}
  \item\label{st:launch} \emph{The launch} \xeqref{foundations}{eq:launch}.
    $\gamma_0 = \val{gamma0}$, $\abs{\vect{p}_0} = \val{p0}$, so
    $\vect{p}_0 = (\val{p0}, 0)$ and $v_0 = \val{v0}$.
  \item\label{st:images} \emph{The fixed images} \xeqref{high_voltage_dome}{eq:kelvin}
    and \eqref{eq:neutral}. The level's charge gives $g$'s first term,
    $\val{imageQUp}$ at $(\val{imageXUp}, \val{imageYUp})$. The reference charge
    gives $\val{imageQPlay}$ at $(\val{imageXPlay}, \val{imageYPlay})$. Their sum
    is $g = \val{inducedG}$, so the centre charge is $Q_c = \val{centreQ}$ at
    $(\val{sphereX}, \val{sphereY})$.
  \item\label{st:march} \emph{The march.} Advance
    \xeqref{high_voltage_dome}{eq:image-energy} and \xeqref{foundations}{eq:eom}
    by \xeqref{foundations}{eq:rk4} with step $\val{H}$, and locate the ending
    face as in \S\ref{sec:flight}. The state every five time units:
    \begin{center}\small\val{states}\end{center}
  \item\label{st:end} \emph{The arrival.} The ion meets the \val{RefFace} face of
    the detector at $(\val{RefX}, \val{RefY})$, at $t = \val{RefT}$, with
    momentum $(\val{RefPx}, \val{RefPy})$, $\val{RefAng}^\circ$ above $+x$.
    The ordinate lies between $\val{yDetLo}$ and $\val{yDetHi}$. The run takes
    $\val{RefSteps}$ steps. The path falls to $\val{RefYmin}$ and then climbs;
    the highest point on the flight is the arrival.
  \item\label{st:check} \emph{The checks.} The fixed potential on the sphere is
    off by at most $\val{surfRes}$. The net charge misses $Q$ by $\val{netQ}$.
    Along the flight, $W$ drifts by at most $\val{RefWdev}$. The step-halving
    estimate of the end state is $\val{RefErr}$. The closest centre distance is
    $\val{RefRclose}$, so the clearance of the surface is $\val{RefClear}$.
  \item\label{st:rad} \emph{The radiation.} Simpson's rule on steps of $h/2$
    gives $\val{RefRad}$, a fraction $\val{RefRadFrac}$ of $T_0$. Halving that
    step again changes the integral by the relative amount $\val{RefRadErr}$.
\end{procedure}

\begin{result}[The reference flight]
The ion enters the detector through its left edge at
$(\val{RefX}, \val{RefY})$ at $t = \val{RefT}$, with momentum
$(\val{RefPx}, \val{RefPy})$. The level is solved.
\end{result}

The march is a few thousand steps of four evaluations. Each evaluation is the
Coulomb sum of seven charges: the two real charges, their two images, the centre
charge, and the ion's image with its following centre charge.

\section{Comparison with the game}\label{sec:comparison}

\comparisontable

Nine comparisons are in the table: the reference outcome, time and end position,
the radiated energy of the reference, and the outcome, time and end position of
the level's own charge with no added charge. The tolerance on a state is
$\val{tolFlight}$. Ten times the largest step-halving estimate $\val{errMax}$
lies above $10^{-9}$, so the tolerance is that product. The radiation tolerance
is $10^{-4}$, relative. The state tolerance and the radiation tolerance were
fixed from these estimates before the comparison.

The largest absolute difference is the reference ordinate, $\val{cmpRefYDiff}$.
The next is the reference time, $\val{cmpRefTDiff}$. The step-halving estimate
$\val{errMax}$ is the same size as the ordinate difference: the fixed step is
what limits the comparison. A verified end state of the engine is finer, about
$10^{-10}$ (\xsecref{foundations}{sec:verification}). Both differences sit under
the tolerance. The bare flight is smoother; its own estimate is $\val{BareErr}$.

The reference abscissa and the bare ordinate agree exactly: each is a face the
flight meets, the detector's left edge and the world's bottom edge. The Simpson
estimate of the radiated energy is $\val{RefRadErr}$. The relative difference
from the game is $\val{cmpRefRadDiff}$, larger than that estimate and far under
the radiation tolerance. The two quadratures are not the same rule; the formula is.

\section{Other solutions}\label{sec:others}

The level's own charge, which is the empty placement in \code{placements.json},
does not solve the level.

The ion still leaves along $+x$, and it never rises above the launch height
$\val{yLaunch}$. The flight leaves through the \val{BareFace} edge of the world
at $(\val{BareX}, \val{BareY})$, at $t = \val{BareT}$, heading $\val{BareAng}^\circ$.
The clearance of the sphere is $\val{BareClear}$, and the energy drifts by
$\val{BareWdev}$. On this placement the images and the centre charge sum to
$\val{bareNet}$.

The positive charge above the launch drives the ion down, and the neutral sphere
does not reverse that. The reference charge is negative and lies just below the
line of flight. The ion is pulled down to $\val{RefYmin}$ and then swung upward,
and it enters the detector still climbing.

\section{What the game computes}\label{sec:engine}

The engine builds the same images (\physref{2.6}). For several spheres it images
again, to a finite depth, and fits a smooth remainder by fundamental solutions
inside each sphere. This level has one sphere, so the second imaging never
starts. A floating sphere then places the centre charge \eqref{eq:neutral}, and
the ion's own pair is \eqref{eq:selfphi}. The surface residual measured in
\S\ref{sec:neutral} is the roundoff $\val{surfRel}$. The energy diagnostic is
\xeqref{high_voltage_dome}{eq:image-energy}. The obstacle is the sphere plus the
contact shell $\val{contact}$.

The flight itself is the adaptive eighth-order Runge--Kutta method of
\physref{5} and \physref{6}, verified as in \physref{7}. The differences in
\S\ref{sec:comparison} match this document's fixed step and sit inside its
tolerance. Radiation is integrated and not applied to the momentum
(\xsecref{foundations}{sec:radiation}, \physref{8}).

\section{Further reading}\label{sec:reading}

The grounded image and the centre charge that raises a sphere to a given
potential are derived for the dome (\xeqref{high_voltage_dome}{eq:kelvin},
\xeqref{high_voltage_dome}{eq:centre}). The centre charge that holds an isolated
sphere at net charge $Q$ is \eqref{eq:neutral}. The flight uses the Foundations'
Runge--Kutta method, Hermite event location and Simpson rule
(\xsecref{foundations}{sec:numerics}). The engine's conductor model is
\physref{2.6}.

\levelfileappendix
\end{document}
